\chapter*{扩写目录与路线图}
\addcontentsline{toc}{chapter}{扩写目录与路线图}

当前书稿已经有一条从电平、逻辑门、状态、ALU、最小 CPU 到整机平台的主线，但它仍然更接近讲义，而不是完整教材。后续扩写不应继续在六章末尾堆专题，而应按经典教材的粒度拆成更多章节：每章有明确问题、结构图、案例、习题和纸面项目。下面是后续扩写采用的目标目录。

\begin{longtable}{L{0.08\linewidth}L{0.25\linewidth}L{0.39\linewidth}L{0.18\linewidth}}
\toprule
章 & 章节名 & 需要讲清楚的问题 & 章末项目 \\
\midrule
1 & 数字抽象与电气基础 & 电压、电流、噪声裕量、阈值、负载、电平标准、受控电流路径如何形成可靠 bit。 & 电平约定图、噪声裕量题 \\
2 & 组合逻辑与布尔代数 & 真值表、布尔代数、NAND/NOR 完备性、译码器、编码器、mux、毛刺和传播延迟。 & 逻辑化简与毛刺案例 \\
3 & 时序逻辑、寄存器和状态机 & 锁存器、触发器、setup/hold、复位、同步器、计数器、移位寄存器和 FSM。 & 状态机分析表 \\
4 & RTL 与硬件描述语言阅读 & 如何读 Verilog/SystemVerilog：组合块、时序块、非阻塞赋值、综合边界、testbench 思路。 & RTL 走读题 \\
5 & 数值表示与算术电路 & 无符号、补码、定点、浮点、BCD、加法器、减法器、乘除法、多精度和标志位。 & 边界向量表 \\
6 & ISA、汇编和机器码 & 指令格式、寄存器、寻址模式、立即数、条件码、调用约定、栈帧和机器码编码。 & 汇编到机器码 trace \\
7 & 数据通路与控制器 & 寄存器堆、ALU、mux、控制字、硬连线控制、微程序控制和数据通路阶段。 & 控制信号表 \\
8 & 单周期与多周期 CPU & IF/ID/EX/MEM/WB 的状态边界，单周期关键路径，多周期复用硬件和等待状态。 & 访存周期图 \\
9 & 流水线、冒险和异常 & 结构冒险、数据冒险、控制冒险、转发、停顿、冲刷、精确异常和中断入口。 & 五级流水线题 \\
10 & cache 与存储层次 & locality、cache line、tag/index/offset、替换、写策略、多级 cache 和一致性入门。 & cache 命中/miss 推演 \\
11 & 虚拟内存、TLB 和页表 & 虚拟地址、页表、TLB、权限、缺页、页错误、IOMMU 和地址空间隔离。 & 地址转换案例 \\
12 & DRAM/DDR 与内存控制器 & SRAM 与 DRAM 差异、bank、row buffer、refresh、ECC、训练、调度和带宽。 & DRAM 访问时序图 \\
13 & 总线、互连和 PCIe & 片选、valid/ready、AXI/APB 概念、PCIe 配置空间、BAR、TLP、MSI 和 AER。 & PCIe 枚举路径 \\
14 & MMIO、中断、DMA 和 IOMMU & 设备寄存器位语义、中断控制器、DMA 描述符、completion、cache 可见性和权限。 & 设备事务 trace \\
15 & SSD、硬盘、网络与外设队列 & 块设备、NVMe、SATA、硬盘、网卡、描述符环、错误恢复、吞吐和尾延迟。 & NVMe/网卡队列题 \\
16 & GPU、VRAM 和显示系统 & GPU 执行模型、SIMD/SIMT、shader、VRAM、命令缓冲、fence、显示扫描和帧缓冲。 & GPU 命令路径图 \\
17 & 主板、固件、启动和平台管理 & VRM、时钟、复位、UEFI/BIOS、PCH、DDR 训练、ACPI、EC、风扇和传感器。 & 启动 bring-up 日志 \\
18 & 链接、装载和系统软件边界 & 目标文件、符号、重定位、ELF、动态链接、loader、系统调用和异常控制流。 & ELF/loader 纸面项目 \\
19 & 性能分析与定量方法 & latency、throughput、CPI、带宽、Amdahl 定律、roofline、cache miss 代价和瓶颈定位。 & 性能分析题 \\
20 & 综合纸面项目与故障案例集 & 从开机、读文件、执行代码、访问设备、GPU 显示到错误恢复的全链路 trace。 & 整机规格与排错报告 \\
\bottomrule
\end{longtable}

\section*{拆分原则}

旧六章不会被简单复制成二十章。拆分时按问题边界移动内容：第一章和第二章保留电气与组合逻辑；当前状态机内容拆到第 3 章；数值和 ALU 拆到第 5、7 章；最小 CPU 拆到第 6、8、9 章；主存、cache、DMA、SSD、PCIe、GPU 和主板分别进入第 10 到第 17 章；目标文件、链接、装载和系统软件边界补成第 18 章；性能分析和综合项目补成第 19、20 章。

当前六章与规划章的对应关系如下，读者可以按右列判断某一主题会先在哪里出现、又会在哪里被完整展开：

\begin{longtable}{L{0.34\linewidth}L{0.56\linewidth}}
\toprule
当前章 & 拆入规划章 \\
\midrule
第 1 章 电平、开关、逻辑门与组合电路 & 第 1、2 章（数字抽象与电气基础、组合逻辑与布尔代数） \\
第 2 章 状态、时钟与同步边界 & 第 3 章（时序逻辑、寄存器和状态机） \\
第 3 章 数值、ALU、数据通路与控制器 & 第 5、7 章（数值表示与算术电路、数据通路与控制器） \\
第 4 章 最小 CPU 与控制 trace & 第 6、8、9 章（ISA 与汇编、单周期与多周期、流水线与冒险） \\
第 5 章 主存、总线、I/O 与 DMA & 第 10–15 章（cache、虚拟内存、DRAM、总线互连、MMIO/DMA、外设队列） \\
第 6 章 经典芯片、启动与平台演进 & 第 16、17 章（GPU 与显示、主板固件与平台管理），并吸收第 13 章 PCIe 内容 \\
第 7 章 动手搭一台 8-bit 教学计算机 & 综合实践章：把第 1–3、6–9 章的规划内容落到可搭建的最小系统 \\
第 8 章 电源、时钟与信号完整性 & 第 1、3 章（电气基础、时序逻辑）的工程延伸，补充电源分配与高速信号内容 \\
第 9 章 RTL 与 Verilog 阅读入门 & 第 4 章（RTL 与硬件描述语言阅读） \\
第 10 章 性能分析与定量方法 & 第 19 章（性能分析与定量方法） \\
第 11 章 外设与队列：SSD、硬盘、网卡 & 第 15 章（SSD、硬盘、网络与外设队列） \\
第 12 章 GPU 与显示系统 & 第 16 章（GPU、VRAM 和显示系统） \\
当前尚未覆盖 & 第 18 章链接装载、第 20 章综合项目 \\
\bottomrule
\end{longtable}

\section*{每章最低质量标准}

每一章至少包含五类内容。第一，开头给出本章要解决的硬件问题，避免堆名词。第二，给出一到三张能说明机制的图，图必须区分数据路径、控制路径、状态边界和错误路径。第三，至少包含一个真实或接近真实的案例，例如 8086 总线、RISC-V 指令、cache miss、NVMe 队列或 GPU fence。第四，章末习题要覆盖概念题、trace 题、边界题和故障定位题。第五，给出参考解答和答题要点，让读者知道答案需要落实到哪些机器细节。

\section*{扩写顺序}

后续优先扩第 6 到第 15 章，因为这些内容决定本书能不能从“硬件入门讲义”变成“计算机底层教材”。建议顺序是：先补 ISA/汇编和数据通路，再补流水线和异常，然后补 cache/TLB/虚拟内存，接着补 DRAM、PCIe、MMIO/DMA、SSD/网络，最后补 GPU、主板启动、链接装载和性能分析。这样每一轮扩写都能增加一个完整知识块，而不是只增加页数。
